\documentclass[a4paper,11pt]{ltjsarticle}
% 数式
\usepackage{amsmath,amsfonts}
\usepackage{bm}
% 画像
\usepackage{graphicx}
\usepackage{circuitikz}
\usepackage{amsmath,amssymb}
\usepackage{siunitx}
\usepackage{float}
\usepackage{tikz}
\usepackage{askmaps}
\usepackage{multirow}
\usepackage{bigstrut}
\usepackage{slashbox}
\usepackage{rotating}
\usepackage{listings}
% 数式
\usepackage{physics}
\usepackage{mathtools}
% 画像
\usepackage{subcaption}
% 表
\usepackage{makecell}
% その他
\usepackage{url}
\usepackage{ascmac}
\usepackage{cases}
\usepackage{here}
\usepackage{upgreek}
% 日本語対応
\usepackage{luatexja}
\usepackage{luatexja-fontspec}

\AtBeginDocument{\RenewCommandCopy\qty\SI}


\definecolor{commentgreen}{RGB}{0,200,0}
\definecolor{eminence}{RGB}{120,80,250}
\definecolor{weborange}{RGB}{255,165,0}
\definecolor{frenchplum}{RGB}{10,150,200}
\definecolor{commentgreen}{RGB}{0,200,0}
\definecolor{eminence}{RGB}{120,80,250}
\definecolor{weborange}{RGB}{255,165,0}
\definecolor{frenchplum}{RGB}{10,150,200}

\lstset{
        language = {C},
        basicstyle = \ttfamily\small,
        keywordstyle=\color{eminence}\ttfamily\bfseries,
        commentstyle=\color{commentgreen}\textit,
    identifierstyle=\color{black}\ttfamily,
        xleftmargin=.35in,
        frame=lines,
    showstringspaces=false,
        numbers=left,
        stepnumber = 1,
        breaklines=true,
        numberstyle = \ttfamily\normalsize,
    tabsize=4,  
        emph={int, int8_t, int16_t, int32_t, int64_t, uint8_t, uint16_t, uint32_t, uint64_t, char, double, float, unsigned, void, bool},
        emphstyle={\color{blue}}, 
        morekeywords={>, <, ., ;, +, -, *, /, !, =, ~},
        breakindent = 10pt, 
        framexleftmargin=10mm, 
        columns=fixed,
        basewidth=0.5em,
        }

% 特定のスタイル設定
\lstdefinestyle{customtxt}{
  basicstyle=\ttfamily\footnotesize,
  backgroundcolor=\color{lightgray},
  frame=single,
  breaklines=true,
  columns=fullflexible,
  showspaces=false,
  showstringspaces=false,
  showtabs=false,
  tabsize=4,
}

\newcommand{\fig}[4]{
    \begin{figure}[htbp]
      \centering
      \includegraphics{./image/#1}
      \caption{#2}
      \label{fig:#3}
    \end{figure}
  }
\begin{document}


\section{要旨}
% 結果等を入れて三行程度でミニレポートを書く
サーミスタは温度を精度よく測定するために作られた半導体であり，半導体は温度を上げていくと電気抵抗が小さくなる．
その特性を温度と電気抵抗を測定することにより確認した．
結果として温度が$4.1$度のときに電気抵抗が$42kΩ$,温度が$77.4$度のときに電気抵抗が$1684Ω$となり，
特性曲線を知ることが出来た．
\section{目的}
% 何を測定して、どのような意味があるのかを書く
サーミスターの電気抵抗と温度計の値を調べて因果関係を調べる．
\section{実験手順}
% 実験指導書のページを参照で良い。しかし違うものをした場合は再現できるほど細かく書く
本実験は追加資料pp.1-4を参考にして行う．ただし温度の測定は$4.1$度から$77.4$度の範囲で15回測定を行った．
測定回数$x$と温度計$y$の値を最小二乗法で求めたところ$y=-5.40x + 82$の関係となり，約$5.4$度ごとの結果を取得している．
また，ホイーストンブリッジの倍率$r_1/r_3$は10回目まで1倍，11回目から15回目まで10倍とした．
そして，実験を始める前にホイートストンブリッジの操作を学ぶために常温での計測を行った．この結果をNo.0とする．
\section{実験結果}
% 表/グラフを用いて実験で得たデータを示し、目的とする物理量を求める。最低限有効数字を意識して、さらには不確かさを表示すること
サーミスターと温度計の測定データを基に表\ref{tab:result}にまとめた．また，グラフを図2方対数グラフを図3に示す．
片対数グラフに引いてある線は最小二乗法により求めた直線である．
\begin{table}[H]
  \centering
  \caption{サーミスターと温度計の測定データ}
  \label{tab:result}
  \begin{tabular}{|c|c|c|c|c|c|}
    \hline
    No. & t(度) & T(K)   & 1/T                 & 1/T-1/T0             & R(Ω)  \\
    \hline
    0   & 25.2 & 298.35 & 3.3518$\cdot 10 ^3$ & -3.0922$\cdot 10 ^4$ & 10340 \\\hline
    1   & 77.4 & 350.55 & 2.8527$\cdot 10 ^3$ & -8.0833$\cdot 10 ^4$ & 1684  \\\hline
    2   & 70.5 & 343.65 & 2.9099$\cdot 10 ^3$ & -7.5105$\cdot 10 ^4$ & 2062  \\\hline
    3   & 64.8 & 337.95 & 2.9590$\cdot 10 ^3$ & -7.0197$\cdot 10 ^4$ & 2470  \\\hline
    4   & 59.8 & 332.95 & 3.0035$\cdot 10 ^3$ & -6.5754$\cdot 10 ^4$ & 2946  \\\hline
    5   & 55.4 & 328.55 & 3.0437$\cdot 10 ^3$ & -6.1732$\cdot 10 ^4$ & 3400  \\\hline
    6   & 50.8 & 323.95 & 3.0869$\cdot 10 ^3$ & -5.7410$\cdot 10 ^4$ & 4000  \\\hline
    7   & 46.4 & 319.55 & 3.1294$\cdot 10 ^3$ & -5.3159$\cdot 10 ^4$ & 4660  \\\hline
    8   & 42.3 & 315.45 & 3.1701$\cdot 10 ^3$ & -4.9092$\cdot 10 ^4$ & 5397  \\\hline
    9   & 35.1 & 308.25 & 3.2441$\cdot 10 ^3$ & -4.1687$\cdot 10 ^4$ & 7005  \\\hline
    10  & 28   & 301.15 & 3.3206$\cdot 10 ^3$ & -3.4039$\cdot 10 ^4$ & 9178  \\\hline
    11  & 20.2 & 293.35 & 3.4089$\cdot 10 ^3$ & -2.5209$\cdot 10 ^4$ & 12070 \\\hline
    12  & 15.4 & 288.55 & 3.4656$\cdot 10 ^3$ & -1.9539$\cdot 10 ^4$ & 15170 \\\hline
    13  & 10.2 & 283.35 & 3.5292$\cdot 10 ^3$ & -1.3179$\cdot 10 ^4$ & 18730 \\\hline
    14  & 6.4  & 279.55 & 3.5772$\cdot 10 ^3$ & -8.3815$\cdot 10 ^5$ & 21850 \\\hline
    15  & 4.1  & 277.25 & 3.6069$\cdot 10 ^3$ & -5.4139$\cdot 10 ^5$ & 24000 \\\hline
  \end{tabular}
\end{table}
\section{考察}
% 実験結果から何がいえるかという問に答えるところ。実験結果が信頼できるのかを論理的に書く。実験指導書の課題もここで答える。新しい実験方法を披露できるとなおよい。最小二乗法は考察
実験結果より，サーミスターの電気抵抗は温度が上がるにつれて小さくなることが分かった．
また，温度と電気抵抗の関係は式\eqref{eq:thermistor}で表されることが知られている．
ただし$T_0$は273.15Kであり，$R_0$は273.15Kのときの電気抵抗である．
それぞれの値を最小二乗法により計算したところ，$R_0 = 29944Ω$，$B = 3.53 \cdot 10^3$，となった．
サーミスターは$R_0$が30kΩになるように作られており，今回の実験結果から得られた値とほぼ一致している．
また，この値を代入することによりホイートストンブリッジの特性曲線を求めることが出来る．
特性曲線を式\eqref{eq:kyokusen}に示す．
\begin{equation}
  R = R_0 e^{ B \cdot \left( \frac{1}{T} - \frac{1}{T_0} \right)}
  \label{eq:thermistor}
\end{equation}
\begin{equation}
  R = 29944 e^{ 3.53 \cdot 10^3 \cdot \left( \frac{1}{T} - \frac{1}{273.15} \right)}
  \label{eq:kyokusen}
\end{equation}




\end{document}